% 6.电路基础.tex 子文件，book文档类，主文档通过\input{6.电路基础.tex}导入
\chapter{第四章 正弦电流电路}

\section{正弦电流}
随正弦规律变化的电流称为正弦电流或正弦交流。

\section{正弦量的相量表示法}
\subsection{复数}

\subsection{正弦量的相量表示}
\[
\dot{A}_m = A_m e^{j\psi} = A_m \angle \psi
\]

\subsection{正弦量运算与相量运算的对应关系}
\begin{enumerate}
    \item 唯一性
    \item 线性性质
    \item 微分规则
\end{enumerate}

\section{基尔霍夫的相量形式}

\section{RLC元件上电压与电流的关系}
\subsection{电阻元件}

\subsection{电感元件}
可得电感上电压、电流相量关系
\[
\]

\subsection{电容元件}

\section{阻抗与导纳}
\subsection{阻抗}
即阻抗端口电压、电流有效值（或振幅）之比等于阻抗模；端口电压超前于电流的相位差等于阻抗角$\psi$。

\subsection{导纳}

\section{正弦电流电路的相量分析法}

\section{正弦电流电路的功率}
\subsection{瞬时功率}

\subsection{平均功率}

\subsection{无功功率}
定义无功功率
\[
\]

\subsection{视在功率}

\subsection{功率因数的提高}

\subsection{复功率}

\section{最大功率传输定理}

\section{耦合电感}
\subsection{互感元件端口特性方程}

\subsection{互感元件的等效电路和去耦合等效}

\section{含互感元件的正弦电流电路}
\subsection{含互感电路的等效化简}

\section{理想变压器}

\endinput